\chapter{Масштаб Хиггса, бегущие константы связи и проблема иерархии}

\lectureoverview{Заключительная лекция начинается с взгляда за пределы Стандартной модели — в сторону объединения, суперсимметрии и проблемы калибровочной иерархии. Чтобы понять, почему поле Хиггса находится в центре этих вопросов, мы сначала восстанавливаем лагранжиан Дирака и показываем, что масса фермиона смешивает левую и правую киральности. Киральное слабое взаимодействие запрещает такую массу как голый член; калибровочно-инвариантная связь Юкавы даёт её только после того, как поле Хиггса приобретает вакуумное значение. Затем мы выводим это вакуумное значение из потенциала Хиггса, спрашиваем, почему каждая известная элементарная масса привязана к этому одному малому масштабу, и используем сверхпроводимость как физическую аналогию. После перерыва лекция выводит ренормализацию из экранирования в проводнике и диэлектрике, переносит рассуждение на QED и QCD, прослеживает калибровочные константы связи в направлении объединения и выводит планковский масштаб из гравитации на малых расстояниях. Затем электрослабый масштаб возвращается как единственное необъяснённое малое число.}

\section{Вопрос, который ждёт нас за пределами Стандартной модели}

Стандартная модель представляется здесь не как замкнутый каталог. В начале сегодняшней лекции я хочу сказать, куда пойдёт следующий курс. У этой теории есть трудности, а не логические парадоксы: проблема калибровочной иерархии, возможная роль суперсимметрии и вопрос об объединении. Три калибровочных множителя
\begin{equation}
  SU(3)_c\times SU(2)_L\times U(1)_Y
  \label{eq:l10-standard-model-group}
\end{equation}
по отдельности описывают сильное, слабое и гиперзарядовое взаимодействия. Могут ли они быть подгруппами одной большей группы? Часто изучаемая возможность — \(SU(5)\), где кварки и лептоны входят в общие мультиплеты. Пользуется ли природа именно этой конструкцией — другой вопрос. Важная подсказка состоит в том, что кажущиеся раздельными структуры удивительно аккуратно укладываются в более крупную математическую структуру.

Большой адронный коллайдер был построен отчасти для того, чтобы выявить физику, которая разрешает эти вопросы. В 2010 его ближайшей целью был бозон Хиггса. Обнаружение Хиггса никак не могло означать просто добавление одной частицы в таблицу. Хиггс опознаётся по взаимосвязанному набору свойств: его вакуумное значение сообщает массы слабым векторным бозонам и фермионам, тогда как его флуктуация взаимодействует с теми же частицами с определёнными силами связи. Резонанс с неправильными связями не был бы хиггсовским бозоном Стандартной модели. \editorialnote{Лекция предшествует открытию 2012 года. Наблюдаемый скаляр с массой приблизительно \(125\,\mathrm{GeV}\) имеет, в пределах нынешней точности, спин, чётность и картину связей, ожидаемые для хиггсовского бозона Стандартной модели. По состоянию на 2026 низкоэнергетических суперсимметричных частиц не наблюдалось. Эти более поздние факты уточняют исторический контекст, не меняя рассуждения, развиваемого в аудитории.}

Путь к этому рассуждению начинается не с потенциала Хиггса, а с уравнения Дирака. Прежде чем спрашивать, откуда берётся масса фермиона, нужно точно понять, что именно она делает.

\section{От уравнения Дирака к его лагранжиану}

\subsection{Старая нотация Дирака}

В гамильтоновой нотации, использовавшейся, когда уравнение Дирака было впервые введено, его структура такова:
\begin{equation}
  i\partial_t\psi
  =
  \left(-i\alpha^i\partial_i+\beta m\right)\psi .
  \label{eq:l10-dirac-hamiltonian}
\end{equation}
Знаки и множители при \(i\) зависят от соглашений, и на доске некоторые из них временно опущены. Инвариантное содержание состоит в разделении между членами, содержащими производные, и массовым членом, не содержащим их.

\begin{figure}[htbp]
\centering
\includegraphics[width=0.82\textwidth]{lecture_10_figure_02.png}
\caption{Уравнение Дирака и лагранжиан в старой нотации \(\alpha,\beta\). Доска используется здесь для показа структуры, а не как фиксированное соглашение о знаках.}
\lectureframe{00:05:52}
\label{fig:l10-dirac-alpha-beta}
\end{figure}

Чтобы получить действие, вариация которого даёт уравнение Дирака, умножим это уравнение на \(\psi^\dagger\). Перенос всех членов в одну сторону устраняет нужду в знаке равенства, что и порождает небольшую математическую шутку этой лекции: как только уравнение записано как выражение, обращающееся в нуль, знак равенства уже сделал своё дело. Соответствующий лагранжиан первого порядка можно схематически записать как
\begin{equation}
  \mathcal L
  =
  \psi^\dagger
  \left(i\partial_t+i\alpha^i\partial_i-\beta m\right)\psi ,
  \label{eq:l10-dirac-old-lagrangian}
\end{equation}
при понимании, что согласованное метрическое соглашение фиксирует относительные знаки.

Члены с производными переносят возбуждение из одной точки пространства-времени в соседнюю и при этом вращают его спинорные компоненты. В решёточной картине это члены перескока. Напротив, массовый член поглощает и вновь испускает фермион в той же самой точке. Производной в нём нет. Вскоре окажется, что эта локальная операция представляет собой повторяющееся превращение между двумя киральностями.

\begin{classroomqa}
\audiencequestion{Нужно ли ставить \(\gamma^0\) после \(\psi^\dagger\)?\lecturetimestamp{00:05:59}}
\lecturerresponse{Не в той нековариантной форме, которая только что была записана. Но вопрос прямо подводит к ковариантной нотации. Определим дираково сопряжение как \(\bar\psi=\psi^\dagger\gamma^0\); тогда \(\gamma^0\) действительно уже входит в объект, стоящий слева.}
\end{classroomqa}

\subsection{Ковариантная форма}

Введём
\begin{equation}
  \bar\psi\equiv\psi^\dagger\beta
  =\psi^\dagger\gamma^0,
  \qquad
  \gamma^0\equiv\beta,
  \qquad
  \gamma^i\equiv\beta\alpha^i .
  \label{eq:l10-dirac-adjoint}
\end{equation}
Поскольку \(\beta^2=1\), умножение на \(\beta\) позволяет свободно переходить между двумя нотациями. Производная по времени принимает вид \(\bar\psi\gamma^0\partial_0\psi\), пространственные производные — вид \(\bar\psi\gamma^i\partial_i\psi\), и в результате получаем
\begin{equation}
  \boxed{\mathcal L_D
  =\bar\psi\left(i\gamma^\mu\partial_\mu-m\right)\psi .}
  \label{eq:l10-covariant-dirac}
\end{equation}

\begin{figure}[htbp]
\centering
\includegraphics[width=0.82\textwidth]{lecture_10_figure_03.png}
\caption{Ковариантная перезапись: \(\bar\psi=\psi^\dagger\gamma^0\), \(\gamma^i=\beta\alpha^i\) и компактный лагранжиан Дирака.}
\lectureframe{00:10:17}
\label{fig:l10-dirac-covariant}
\end{figure}

Эта нотация делает явной лоренцеву природу каждой билинейной комбинации. \(\bar\psi\psi\) — скаляр, тогда как \(\bar\psi\gamma^\mu\psi\) — четырёхвектор. Компактное выражение не просто красивее формы \(\alpha,\beta\); оно показывает, как преобразуются его части.

\begin{classroomqa}
\audiencequestion{Должен ли был спинор в самом левом месте компактного выражения быть \(\bar\psi\)?\lecturetimestamp{00:09:39}}
\lecturerresponse{Да. Пропущенная черта восстановлена на доске. Это важно: \(\bar\psi\) — не \(\psi^\dagger\), а \(\psi^\dagger\gamma^0\).}
\end{classroomqa}

\section{Киральность и то, что на самом деле делает дираковская масса}

Помимо \(\gamma^0,\gamma^1,\gamma^2,\gamma^3\), введём
\begin{equation}
  \gamma^5
  =i\gamma^0\gamma^1\gamma^2\gamma^3,
  \qquad
  (\gamma^5)^2=1 .
  \label{eq:l10-gamma-five}
\end{equation}
Показанный здесь множитель \(i\) соответствует стандартному соглашению для метрики Минковского с преимущественно отрицательными знаками. Собственные значения равны \(+1\) и \(-1\), а соответствующие проекторы имеют вид
\begin{equation}
\begin{aligned}
  P_L&=\frac{1-\gamma^5}{2},
  &
  P_R&=\frac{1+\gamma^5}{2},
  \\
  \psi_L&=P_L\psi,
  &
  \psi_R&=P_R\psi .
\end{aligned}
  \label{eq:l10-chiral-projectors}
\end{equation}

\begin{classroomqa}
\audiencequestion{Почему дополнительная матрица называется \(\gamma^5\), если \(\gamma^4\) нет? И описывает ли её собственное значение спиральность?\lecturetimestamp{00:10:42}}
\lecturerresponse{Название историческое: это дополнительная матрица, образованная как произведение четырёх пространственно-временных гамма-матриц. Её собственное значение помечает киральность, а не, вообще говоря, спиральность. Для безмассовой частицы эти два понятия при надлежащем отождествлении можно считать совпадающими, поэтому в языке их легко смешать, но для слабой калибровочной теории нужна именно киральность как метка лоренцева представления.}
\end{classroomqa}

Кинетический член сохраняет киральность. Один удобный способ это увидеть состоит в том, что \(\{\gamma^\mu,\gamma^5\}=0\), так что он распадается на
\begin{equation}
  \bar\psi i\gamma^\mu\partial_\mu\psi
  =
  \bar\psi_L i\gamma^\mu\partial_\mu\psi_L
  +
  \bar\psi_R i\gamma^\mu\partial_\mu\psi_R .
  \label{eq:l10-chiral-kinetic}
\end{equation}
Скалярная билинейная комбинация ведёт себя иначе:
\begin{equation}
  \bar\psi\psi
  =
  \bar\psi_L\psi_R+\bar\psi_R\psi_L .
  \label{eq:l10-chiral-mass}
\end{equation}

\begin{figure}[htbp]
\centering
\includegraphics[width=0.80\textwidth]{lecture_10_figure_04.png}
\caption{На доске скалярная билинейная комбинация разложена на лево-правую и право-левую части.}
\lectureframe{00:15:38}
\label{fig:l10-chiral-decomposition}
\end{figure}

Таким образом, дираковская масса — это локальная операция, меняющая киральность:
\begin{equation}
  \psi_L\longleftrightarrow\psi_R .
\end{equation}
Безмассовый фермион может распространяться с определённой киральностью. Массивный же не может оставаться в одном киральном секторе, потому что массовый член непрерывно смешивает оба. Может показаться удивительным, что это превращение вообще имеет отношение к инерции, но именно так в релятивистской теории поля выглядит масса покоя фермиона.

\section{Почему слабая симметрия запрещает голую массу фермиона}

Заряженный слабый ток имеет чисто левокиральную структуру. На лекции сначала называется \(Z\), а затем формулировка исправляется на \(W^\pm\): связи \(Z\) киральны и неодинаковы, но они не просто отсутствуют для любого правокирального фермиона. Корректная формулировка такова:
\begin{equation}
  \mathcal L_{\mathrm{cc}}
  =
  -\frac{g}{\sqrt2}
  \left(\bar u_L\gamma^\mu d_L
  +\bar\nu_L\gamma^\mu e_L\right)W^+_\mu
  +\mathrm{h.c.}
  \label{eq:l10-charged-current}
\end{equation}
Левокиральные фермионы преобразуются как дублеты \(SU(2)_L\); правокиральные заряженные фермионы преобразуются как синглеты.

Теперь сравним эти правила преобразования с \eqref{eq:l10-chiral-mass}. Член вида \(\bar e_L e_R\) поглощает объект с одним набором электрослабых квантовых чисел и испускает объект с другим. Сам по себе он не является калибровочно-инвариантным. Поэтому привычная дираковская масса недопустима как голый член в электрослабом лагранжиане.

Именно таков правильный порядок объяснения. Массы фермионов не вводятся сначала, а затем чисто косметически связываются с полем Хиггса. Калибровочная симметрия сначала устраняет очевидный массовый член. Затем поле Хиггса вынуждено нести недостающие квантовые числа.

\subsection{Поле Хиггса как носитель заряда}

Пусть поле Хиггса преобразуется как дублет \(SU(2)_L\) с гиперзарядом \(+\tfrac12\):
\begin{equation}
  H=
  \begin{pmatrix}
    H^+\\ H^0
  \end{pmatrix},
  \qquad
  \widetilde H=i\sigma^2H^* .
  \label{eq:l10-higgs-doublet}
\end{equation}
На лекции групповые индексы оставлены схематическими, а поле обозначено как \(\phi\). В этой нотации член Юкавы имеет вид
\begin{equation}
  \mathcal L_Y
  =
  -G_Y\left(
    \bar\psi_L\phi\psi_R+
    \bar\psi_R\phi^\dagger\psi_L
  \right).
  \label{eq:l10-yukawa-schematic}
\end{equation}
Поле Хиггса уносит слабый заряд, который голое изменение киральности не сохранило бы.

Полное выражение Стандартной модели для одного поколения делает тот же учёт явным:
\begin{equation}
  \mathcal L_Y
  =
  -y_d\bar Q_LHd_R
  -y_u\bar Q_L\widetilde H u_R
  -y_e\bar L_LHe_R
  +\mathrm{h.c.}
  \label{eq:l10-yukawa-full}
\end{equation}
Каждый член калибровочно-инвариантен, и каждый вид фермиона имеет свою собственную безразмерную константу связи Юкавы.

\begin{figure}[htbp]
\centering
\includegraphics[width=0.80\textwidth]{lecture_10_figure_05.png}
\caption{Переворот киральности с участием Хиггса и потенциал нарушения симметрии показаны на доске.}
\lectureframe{00:18:54}
\label{fig:l10-yukawa-potential-board}
\end{figure}

\section{Одно взаимодействие даёт и массу, и взаимодействие с Хиггсом}

Природа выбирает вакуум, в котором нейтральная компонента Хиггса не обращается в нуль. В схематической нотации лекции для вещественного поля
\begin{equation}
  \phi=F+\mathcal H ,
  \label{eq:l10-f-plus-h}
\end{equation}
где \(F\) — нефлуктуирующий вакуумный сдвиг, а \(\mathcal H\) — распространяющаяся флуктуация Хиггса. Подстановка в
\eqref{eq:l10-yukawa-schematic} даёт
\begin{align}
  \mathcal L_Y
  ={}&
  -G_YF
  \left(\bar\psi_L\psi_R+\bar\psi_R\psi_L\right)
  \notag\\
  &-G_Y\mathcal H
  \left(\bar\psi_L\psi_R+\bar\psi_R\psi_L\right).
  \label{eq:l10-yukawa-split}
\end{align}
Первая строка — это масса Дирака, а вторая — взаимодействие с физическим квантом Хиггса. В нормировке, принятой в лекции,
\begin{equation}
  m_f=G_YF
  \qquad\hbox{in the lecture's normalization.}
  \label{eq:l10-mass-lecture-normalization}
\end{equation}

В стандартной нормировке Стандартной модели
\begin{equation}
\begin{aligned}
  H(x)&=\frac{1}{\sqrt2}
  \begin{pmatrix}
    0\\ v+h(x)
  \end{pmatrix},
  \\
  m_f&=\frac{y_fv}{\sqrt2},
  \\
  \mathcal L_{hff}&=-\frac{m_f}{v}h\bar f f,
\end{aligned}
  \label{eq:l10-standard-higgs-normalization}
\end{equation}
где \(v\simeq246\,\mathrm{GeV}\). Приближённая
\(F\sim200\,\mathrm{GeV}\) в лекции — это тот же самый электрослабый масштаб с опущенными деталями нормировки.

Измеренные массы \(W\) и \(Z\) определяют масштаб \(F\). Измеренные массы фермионов затем определяют юкавские константы связи. Константа связи электрона крошечна; у топ-кварка она порядка единицы. Если бы этим действие механизма исчерпывалось, мы бы вводили по одному неизвестному числу на каждый фермион и затем извлекали бы ровно ту же информацию обратно.

Но вторая строка \eqref{eq:l10-yukawa-split} несёт новое экспериментальное содержание. Она говорит, что фермион может испустить или поглотить бозон Хиггса и, если рассматривать скрещённый канал, что бозон Хиггса может распадаться на пару фермион-антифермион. Схематически,
\begin{equation}
\begin{aligned}
  \mathcal M(h\to f\bar f)&\propto y_f,
  \\
  \Gamma(h\to f\bar f)
  &=
  N_c\frac{y_f^2m_h}{16\pi}
  \\
  &\quad\times
  \left(1-\frac{4m_f^2}{m_h^2}\right)^{3/2},
\end{aligned}
  \label{eq:l10-higgs-fermion-width}
\end{equation}
на древесном уровне. Здесь \(N_c=3\) для кварков и \(1\) для лептонов. Большая масса фермиона означает большую юкавскую константу связи, если только распад кинематически разрешён. Это не означает, что любой тяжёлый канал доминирует: важны также цветовые множители, фазовое пространство, каналы, индуцированные петлями, и распады на векторные бозоны. Это означает, что сама константа связи с фермионами пропорциональна массе.

\begin{classroomqa}
\audiencequestion{Как на практике предсказать время жизни бозона Хиггса, не вдаваясь во все технические детали?\lecturetimestamp{00:25:32}}
\lecturerresponse{Взаимодействие задаёт матричный элемент гамильтониана. Нарисуйте соответствующую диаграмму Фейнмана, вычислите соответствующую амплитуду перехода, возведите амплитуду в квадрат и просуммируйте или проинтегрируйте по всем допустимым конечным состояниям. Результат — скорость распада. Время жизни обратно полной ширине, \(\tau=\hbar/\Gamma_{\mathrm{tot}}\). Это та же самая квантовомеханическая логика, которая используется в расчётах частиц с раннего этапа развития квантовой теории поля.}
\end{classroomqa}

На момент этой лекции такая феноменология Хиггса всё ещё была предсказанием, а не прямым наблюдением. Самовзаимодействия \(W\) и \(Z\) уже дали жёсткие проверки калибровочной структуры; бозон Хиггса и его прямые взаимодействия ещё предстояло обнаружить. Эта историческая неопределённость должна остаться в изложении. Как и более поздний итог: была обнаружена частица, подобная Хиггсу, и измерения её взаимодействий теперь служат проверками той же самой структуры порождения масс, которая здесь развивается.

\section{Потенциал, определяющий F}

Вакуумный сдвиг не выбирают вручную после записи юкавской связи. Он определяется потенциалом Хиггса. Для одного вещественного среза полного дублета возьмём
\begin{equation}
  V(\phi)
  =
  -\frac{\mu^2}{2}\phi^2
  +\frac{\lambda}{4}\phi^4,
  \qquad
  \mu^2>0,\quad\lambda>0.
  \label{eq:l10-real-higgs-potential}
\end{equation}
Вблизи \(\phi=0\) отрицательный квадратичный член даёт перевёрнутую параболу. Называть \(\mu\) мнимой массой — это сокращённый способ сказать, что начало координат неустойчиво; это не означает, что физическое возбуждение Хиггса около истинного вакуума имеет мнимую массу.

Один только квадратичный член уводил бы поле без границ и не оставлял бы состояния с наименьшей энергией. Положительный член четвёртой степени в конце концов начинает доминировать и загибает потенциал вверх. Нечётная степень нарушила бы игрушечную симметрию \(\phi\mapsto-\phi\). Член четвёртой степени — это первое чётное стабилизирующее взаимодействие и, в четырёх измерениях пространства-времени, скалярное взаимодействие наибольшей полиномиальной степени с безразмерной константой связи.

Продифференцируем:
\begin{align}
  \frac{dV}{d\phi}
  &=-\mu^2\phi+\lambda\phi^3
  \notag\\
  &=\phi\left(-\mu^2+\lambda\phi^2\right).
  \label{eq:l10-potential-derivative}
\end{align}
Стационарная точка \(\phi=0\) — это неустойчивый максимум. Два минимума удовлетворяют
\begin{equation}
  \boxed{F^2=\frac{\mu^2}{\lambda}},
  \qquad
  F=\frac{\mu}{\sqrt\lambda}
  \label{eq:l10-vacuum-value}
\end{equation}
в этой конвенции с вещественным полем.

\begin{figure}[htbp]
\centering
\includegraphics[width=0.84\textwidth]{lecture_10_figure_06.png}
\caption{Завершённая минимизация на доске:
\(V=-\mu^2\phi^2/2+\lambda\phi^4/4\),
\(V'=0\) и \(\phi^2=\mu^2/\lambda\).}
\lectureframe{00:34:32}
\label{fig:l10-higgs-minimization}
\end{figure}

Множители \(1/2\) и \(1/4\) были введены именно затем, чтобы они исчезли при дифференцировании. Вблизи начала координат \(\phi^4\) меньше, чем \(\phi^2\); при больших значениях поля этот член растёт быстрее и стабилизирует систему. Такое простое соперничество порождает двойную яму.

Размерности также проясняют мимолётную словесную оговорку в расшифровке. \(\lambda\) безразмерна в четырёх измерениях пространства-времени, тогда как \(\mu\) и \(F\) имеют размерность массы. Разлагая \(\phi=F+\mathcal H\), из кривизны потенциала получаем
\begin{equation}
  m_{\mathcal H}^2
  =
  \left.\frac{d^2V}{d\phi^2}\right|_{\phi=F}
  =
  -\mu^2+3\lambda F^2
  =
  2\mu^2
  =
  2\lambda F^2.
  \label{eq:l10-higgs-curvature}
\end{equation}
В полной записи Стандартной модели,
\(V(H)=-\mu^2H^\dagger H+\lambda(H^\dagger H)^2\), аналогично имеем
\(m_h^2=2\lambda v^2\). Измеренные масса Хиггса и \(v\) дают
\(\lambda\) около \(0.13\) на электрослабом масштабе. Оценка лекции 2010 года была намеренно грубой; структурное утверждение состояло в том, что \(\lambda\) не является абсурдно малым числом.

\section{Почему всё, что мы наблюдаем, связано именно с этим масштабом?}

Один и тот же \(F\) определяет массы \(W\) и \(Z\) и, будучи умноженным на разные юкавские константы связи, также все массы заряженных фермионов. Такое эмпирическое сосредоточение ставит вопрос, более интересный, чем сама алгебра: почему все известные элементарные частицы должны заимствовать свои массы из одного и того же явления нарушения симметрии?

Нет теоремы, что так должно быть для любой возможной частицы. Представьте вектороподобный фермион, у которого левая и правая киральные компоненты преобразуются одинаково относительно калибровочной группы. Тогда
\(\bar\psi_L\psi_R+\mathrm{h.c.}\) может быть калибровочно-инвариантным и без Хиггса. Полный калибровочный синглет точно так же может иметь массу, не связанную с электрослабым нарушением. Такие частицы естественным образом могли бы быть гораздо тяжелее частиц, привязанных к \(F\). Обнаруженные нами частицы могут быть достаточно лёгкими, чтобы их можно было увидеть, именно потому, что симметрия делала бы их безмассовыми, когда \(F\) обращается в нуль.

Следовательно, исключительным числом является электрослабый масштаб. Спекулятивная картина такова: типичные незащищённые массы лежат вблизи гораздо большего фундаментального масштаба, тогда как частицы, защищённые киральной калибровочной симметрией, утягиваются вниз к \(F\).

\begin{classroomqa}
\audiencequestion{Является ли планковская масса одним из кандидатов на фундаментальный масштаб масс, с которым следует сравнивать \(F\)?\lecturetimestamp{00:38:43}}
\lecturerresponse{Да. Другой кандидат — масштаб, на который указывает унификация калибровочных констант связи. Оба возникнут из доводов далее в лекции, хотя никакие частицы с такими огромными массами непосредственно не наблюдались.}
\end{classroomqa}

\begin{classroomqa}
\audiencequestion{Что значит утверждение, что другие масштабы примерно на 15 порядков величины выше \(F\), и каким частицам они соответствуют?\lecturetimestamp{00:39:54}}
\lecturerresponse{Электрослабый масштаб имеет порядок
\(10^2\,\mathrm{GeV}\), тогда как унификация предполагается вблизи
\(10^{16}\,\mathrm{GeV}\), а планковский масштаб лежит вблизи
\(10^{19}\,\mathrm{GeV}\). Мы не создавали частицы таких масс.
О масштабах мы судим косвенно: об одном — экстраполируя измеренные калибровочные константы связи, о другом — комбинируя квантовую механику, теорию относительности и ньютоновскую постоянную.}
\end{classroomqa}

\section{Заряженный конденсат в лаборатории}

Перед перерывом из аудитории задают вопрос об обычном аналоге механизма Хиггса в физике конденсированного состояния. Ответ развивается поэтапно, а не сводится к одной лишь аналогии.

\begin{classroomqa}
\audiencequestion{Помимо ультрахолодных бозе-конденсатов, есть ли распространённый пример хиггсоподобного механизма в физике твёрдого тела?\lecturetimestamp{00:40:50}}
\lecturerresponse{Да. Нейтральное бозонное поле может приобретать ненулевое среднее значение в бозе-конденсате или сверхтекучей жидкости. Если конденсирующийся объект заряжен, то то же самое основное явление изменяет калибровочное поле. Знакомый пример даёт сверхпроводимость: заряженные куперовские пары конденсируются, и электромагнитные поля приобретают внутри материала конечную глубину проникновения.}
\end{classroomqa}

На молекулярном уровне разрешения атом гелия можно рассматривать как одну эффективную бозонную частицу с полем \(\Phi_{\mathrm{He}}\). Пустому пространству соответствует нулевое поле; в сверхтекучей фазе имеет место
\(\langle\Phi_{\mathrm{He}}\rangle\ne0\), и соответствующая глобальная фазовая симметрия спонтанно нарушается. Атомы гелия нейтральны, так что это ещё не механизм массы калибровочного поля.

В сверхпроводнике электроны образуют слабо связанные куперовские пары. Поле пары схематически представляет собой билинейную комбинацию электронных полей,
\begin{equation}
  \Delta(x)\sim\langle\psi_\downarrow(x)\psi_\uparrow(x)\rangle,
  \label{eq:l10-cooper-order-parameter}
\end{equation}
и несёт электрический заряд \(2e\). Когда \(\Delta\) конденсируется, электромагнитное калибровочное поле реагирует ровно так, как и должно реагировать калибровочное поле, связанное с заряженным параметром порядка. Статические магнитные поля затухают на лондоновской глубине проникновения:
\begin{equation}
  B(x)\propto e^{-x/\lambda_L}.
  \label{eq:l10-meissner-screening}
\end{equation}
Обратная глубина проникновения играет роль эффективной массы фотона внутри среды. Это не масса фотона в вакууме; это свойство сверхпроводящей фазы.

\begin{classroomqa}
\audiencequestion{Разве два электрона в куперовской паре не разделены,
даже если они лежат по разные стороны поверхности
Ферми?\lecturetimestamp{00:43:13}}
\lecturerresponse{Это не крошечная молекула, локализованная в одной точке. Куперовская
пара является протяжённой и естественным образом описывается в пространстве
импульсов коррелированными состояниями вблизи противоположных импульсов. Однако
для аргумента о симметрии она ведёт себя как единый заряженный бозонный параметр
порядка. Важно именно её конденсирование.}
\end{classroomqa}

Исторически эта связь — не просто учебная метафора.
Сверхпроводимость и анализ Андерсоном калибровочных полей в фазе с нарушенной
симметрией помогли сделать релятивистский механизм Хиггса понятным. Эта
конструкция в физике частиц была разработана независимо в релятивистской
теории поля, но общая математическая структура здесь действительно одна и та же.

\section{Голые параметры и измеряемые параметры}

После перерыва мы возвращаемся к лагранжиану. Он содержит
параметры, называемые массами и зарядами, но эти входные параметры редко
тождественны величинам, измеряемым без дополнительных оговорок. Высокочастотные
моды взаимодействуют с низкочастотными модами и образуют их одетые состояния.
Систематическая связь между параметрами, определёнными при одном разрешении, и
наблюдаемыми при другом, — это ренормализация.

Электрический заряд позволяет построить наглядную картину без вычислений.
На больших расстояниях наблюдаемый заряд определяется законом Кулона:
\begin{equation}
  V(r)\sim\frac{e_1e_2}{4\pi r},
  \qquad
  F(r)\sim\frac{e_1e_2}{4\pi r^2}.
  \label{eq:l10-coulomb-law}
\end{equation}
В лекции множитель \(4\pi\) опускается, и пишут \(e^2/r\) или \(e^2/r^2\).
Существенно то, что коэффициент, выводимый из измерения силы, может
зависеть от расстояния, на котором исследуется система.

\subsection{Проводник: полное экранирование}

Поместите положительный пробный заряд в идеальный проводник. Его электрическое поле
заставляет подвижный отрицательный заряд смещаться к нему, а положительный — от
него. Ток течёт до тех пор, пока локальный источник не окажется окружён облаком
противоположного заряда и в толще вещества вне этого облака не останется поля;
компенсирующий заряд вытесняется к удалённой границе. Отрицательный источник
вызывает противоположную перестройку. Тогда два одетых источника, находящиеся
далеко друг от друга, могут практически не действовать друг на друга, хотя
микроскопически внесённые заряды были ненулевыми.

\begin{figure}[htbp]
\centering
\includegraphics[width=0.82\textwidth]{lecture_10_figure_07.png}
\caption{Заряд источника, окружённый экранирующим облаком противоположного
заряда, в обсуждении проводника.}
\lectureframe{00:50:24}
\label{fig:l10-conductor-screening}
\end{figure}

\begin{classroomqa}
\audiencequestion{Именно эта перестройка и имеется в виду под
экранированием?\lecturetimestamp{00:51:25}}
\lecturerresponse{Да. На расстояниях, больших по сравнению с размером облака, внесённый
заряд экранирован. Если ввести другой зонд внутрь облака, часть
экранирующего облака уже не может находиться между источником и зондом;
тогда исходный заряд проявляется сильнее.}
\end{classroomqa}

Облако имеет характерный масштаб, задаваемый динамикой и плотностью
подвижных зарядов. При расстояниях \(r\), меньших этого масштаба, зонд
проникает внутрь облака. Поэтому можно определить бегущий заряд \(e(r)\)
через
\begin{equation}
  F(r)\equiv\frac{e^2(r)}{4\pi r^2}.
  \label{eq:l10-running-charge-definition}
\end{equation}
Для идеализированного проводника \(e(r)\) стремится к нулю на очень
больших расстояниях и к микроскопическому значению на очень малых.

\subsection{Диэлектрик: частичное экранирование}

В изоляторе электроны связаны с атомами, а не свободно перемещаются через
образец. Грубая, но полезная модель представляет каждое электронное облако
прикреплённым к своему иону пружиной. Без внешнего поля средние центры
положительного и отрицательного зарядов совпадают. Приложенное поле слегка
смещает их друг относительно друга и индуцирует диполь.

Вокруг отрицательного источника близлежащие электронные облака смещаются от него,
оставляя свои положительные ионы немного ближе; вокруг положительного источника
они смещаются в другую сторону. Связанные заряды не могут сместиться достаточно
далеко, чтобы полностью нейтрализовать источник, поэтому поле на больших
расстояниях уменьшается благодаря диэлектрической проницаемости, а не
устраняется:
\begin{equation}
  V_{\mathrm{medium}}(r)
  =
  \frac{1}{\epsilon_r}
  \frac{e_1e_2}{4\pi\epsilon_0r}.
  \label{eq:l10-dielectric-coulomb}
\end{equation}
На расстояниях внутри облака поляризации зонд снова видит большую часть
голого источника. Полное и частичное экранирование — две версии одной и той же
масштабно-зависимой идеи.

\section{Вакуум КЭД как поляризуемая среда}

Вакуум квантовой электродинамики не содержит материальных атомов, но в нём есть
заряженные квантовые поля. Флуктуации электрон-позитронных пар дают петлевые
вклады в распространение фотона. В старой терминологии моря Дирака электрическое
поле искажает заполненные состояния с отрицательной энергией; в современной
терминологии поляризация вакуума одевает фотонный пропагатор. Согласно картине
среды, виртуальная пара поляризуется: её положительная и отрицательная
компоненты смещаются в противоположных направлениях.

Положительный источник окружён небольшим избытком отрицательного
поляризационного заряда, а отрицательный — положительным поляризационным зарядом.
Поэтому заряд, измеряемый издалека, меньше параметра, который связали бы с
неразрешённым источником. В отличие от обычного вещества,
вакуум содержит флуктуации в непрерывном диапазоне импульсов. Низкоэнергетические
пары соответствуют структуре на больших расстояниях; высокоэнергетические пары
зондируют меньшие расстояния. Не существует конечного радиуса облака, на котором
этот процесс просто прекращается.

Когда два электрона подносят ближе друг к другу, сначала перестают помещаться
между ними самые большие облака поляризации. При дальнейшем сближении
последовательно снимаются всё меньшие облака. Эффективная электрическая
константа связи растёт логарифмически с обратным расстоянием, или,
что эквивалентно, с энергией:
\begin{equation}
  \alpha(Q)
  \simeq
  \frac{\alpha(\mu)}
  {1-\dfrac{2\alpha(\mu)}{3\pi}
  \ln\!\left(\dfrac{Q}{\mu}\right)}
  \label{eq:l10-qed-running}
\end{equation}
для одного заряженного дираковского фермиона далеко выше его порога в ведущем
порядке. Дополнительные виды заряженных частиц меняют коэффициент при
прохождении их порогов.

\begin{figure}[htbp]
\centering
\includegraphics[width=0.82\textwidth]{lecture_10_figure_08.png}
\caption{Эффективный заряд КЭД медленно возрастает по мере увеличения
горизонтального масштаба \(1/R\) и, следовательно, энергии зондирования.}
\lectureframe{01:02:13}
\label{fig:l10-qed-running-board}
\end{figure}

Электронное рассеяние даёт операциональную формулировку этого. Более высокая
энергия столкновения означает меньшую длину волны де Бройля и меньшее расстояние
наибольшего сближения. Поэтому в амплитуде должна использоваться константа связи,
соответствующая данной передаче импульса, а не одно фиксированное число для всех
масштабов.

\begin{classroomqa}
\audiencequestion{На каком расстоянии эта зависимость от масштаба впервые становится
измеримой?\lecturetimestamp{01:03:54}}
\lecturerresponse{Поляризация вакуума электронами становится существенной вблизи
комптоновского масштаба электрона, \(\lambda_C=\hbar/(m_ec)\), примерно
\(4\times10^{-11}\,\mathrm{cm}\) для приведённой длины волны. Эффект
мал и логарифмичен, поэтому требуется высокая точность. Более тяжёлые заряженные
частицы подключаются к бегу на своих собственных порогах.}
\end{classroomqa}

\begin{classroomqa}
\audiencequestion{Не предполагает ли разговор о наименьшем расстоянии, что мы
уже знаем, что там происходит?\lecturetimestamp{01:04:49}}
\lecturerresponse{Нет. Мы экстраполируем теорию лишь до тех пор, пока её состав
частиц и взаимодействия остаются применимыми. Новая физика может изменить бег.
Эксперимент не говорит нам, что КЭД остаётся полным описанием на сколь угодно
малых расстояниях.}
\end{classroomqa}

\begin{classroomqa}
\audiencequestion{Для глюонов зависимость от масштаба противоположна? Почему?
\lecturetimestamp{01:05:03}}
\lecturerresponse{Да. Кварковые петли экранируют цветовой заряд по аналогии
с заряженной материей, но сами глюоны несут цвет и нелинейно взаимодействуют.
Их вклад имеет противоположный знак и доминирует при наблюдаемом числе кварковых
ароматов. Результат — антиэкранирование и асимптотическая свобода.}
\end{classroomqa}

\begin{classroomqa}
\audiencequestion{Как можно было бы определить голый электрический заряд, если для этого
требуется бесконечно малое расстояние?\lecturetimestamp{01:05:28}}
\lecturerresponse{Нельзя непосредственно измерить единственный, не зависящий от
регуляризации, голый заряд. Параметр определяют при выбранной отсечке по малым
расстояниям или на выбранном масштабе ренормализации, а затем подстраивают его
так, чтобы некоторая измеренная низкоэнергетическая величина оставалась
фиксированной. Масштабно-зависимая ренормированная константа связи и есть
полезный физический объект.}
\end{classroomqa}

\begin{classroomqa}
\audiencequestion{Стремится ли заряд КЭД к конечному значению при бесконечной
энергии?\lecturetimestamp{01:05:39}}
\lecturerresponse{Нет, если экстраполировать одну лишь пертурбативную КЭД. Её
положительная бета-функция заставляет константу связи продолжать расти
логарифмически и формально приводит к полюсу Ландау на фантастически высоком
масштабе. Задолго до того, как этот формальный полюс пришлось бы трактовать как
физический, ожидается, что описание будет вложено в электрослабую теорию и в
конечном счёте изменено дальнейшей физикой.}
\end{classroomqa}

\subsection{Что ренормализация изменяет и что сохраняет}

Выберите длину отсечки \(a\), обозначьте соответствующий входной параметр как \(e(a)\) и
вычислите некоторую фиксированную наблюдаемую на больших расстояниях. Другая
отсечка \(a'\) столь же допустима, если входной параметр изменён на \(e(a')\) так, чтобы
воспроизводилась та же наблюдаемая:
\begin{equation}
  \frac{d}{d\ln\mu}\,\mathcal O_{\mathrm{physical}}=0,
  \qquad
  \mu\frac{de}{d\mu}=\beta(e).
  \label{eq:l10-rg-invariance}
\end{equation}
Этот компенсирующий поток и есть ренормгруппа.

\begin{classroomqa}
\audiencequestion{Должно ли слово «голый» относиться именно к планковской
длине?\lecturetimestamp{01:06:20}}
\lecturerresponse{Нет. Можно остановить описание на любой выбранной
отсечке и назвать параметр там голым. Изменение отсечки при соответствующем
изменении этого параметра оставляет доступную физику неизменной. Планковская
длина — правдоподобное место, где обычная квантовая теория поля должна быть
заменена, а не обязательная конвенция ренормализации.}
\end{classroomqa}

\begin{classroomqa}
\audiencequestion{Является ли эта процедура, по сути, приближением и огрублением описания?\lecturetimestamp{01:07:15}}
\lecturerresponse{Это концептуальное приближение. Мы спрашиваем, как должны изменяться параметры более детального описания, чтобы предсказания при фиксированной, экспериментально доступной длине волны оставались теми же.}
\end{classroomqa}

\begin{classroomqa}
\audiencequestion{Наблюдается ли у гравитации то же явление бега?\lecturetimestamp{01:08:14}}
\lecturerresponse{Гравитация тоже зависит от масштаба, но не просто через тот же самый механизм диэлектрического экранирования. Для нынешней лекции достаточно размерностного аргумента: локализация повышает энергию, энергия гравитирует, и возникающая безразмерная сила гравитационного взаимодействия быстро растёт с энергией.}
\end{classroomqa}

\section{КХД: антиэкранирование, трубки потока и асимптотическая свобода}

На адронных расстояниях силовые линии хромоэлектрического поля не расходятся сферически, как в электромагнитном кулоновском поле. Самодействия глюонов сжимают их в узкую трубку между кварком и антикварком. Если у этой трубки энергия на единицу длины примерно постоянна, \(\sigma\), то
\begin{equation}
  V_{q\bar q}(r)
  \simeq
  -\frac{C_F\alpha_s(r)}{r}
  +\sigma r .
  \label{eq:l10-cornell-potential}
\end{equation}
При больших \(r\) линейный член доминирует. Растягивание пары создаёт дополнительный участок трубки, как при вытягивании тягучей ириски, и энергия растёт, а не стремится к нулю. Это сторона этой истории, связанная с конфайнментом.

\begin{figure}[htbp]
\centering
\includegraphics[width=0.84\textwidth]{lecture_10_figure_09.png}
\caption{Силовые линии хромоэлектрического поля изображаются как трубка, соединяющая разделённые цветовые источники, а не как свободно расходящееся кулоновское поле.}
\lectureframe{01:09:56}
\label{fig:l10-qcd-flux-tube-board}
\end{figure}

На расстояниях, меньших радиуса трубки, потенциал приблизительно возвращается к кулоновской форме, но его константа связи убывает с энергией. В однопетлевом приближении,
\begin{equation}
  \alpha_s(Q)
  =
  \frac{4\pi}
  {\left(11-\frac23n_f\right)
  \ln(Q^2/\Lambda_{\mathrm{QCD}}^2)} .
  \label{eq:l10-asymptotic-freedom}
\end{equation}
Для \(n_f<17\) коэффициент положителен, и \(\alpha_s(Q)\to0\) при \(Q\to\infty\). КХД слаба на очень малых расстояниях и сильна на больших. Это точная версия противоположного поведения, обсуждавшегося на лекции, в терминах бегущей константы связи.

\begin{classroomqa}
\audiencequestion{Что происходит со слабой константой связи между кривыми КЭД и КХД?\lecturetimestamp{01:11:53}}
\lecturerresponse{Её бег занимает промежуточное положение и зависит от того, какой электрослабый множитель отслеживается. Выше электрослабого масштаба правильными величинами являются константы связи \(U(1)_Y\) и \(SU(2)_L\). Одна растёт, а другая медленно убывает. Вместе с убывающей сильной константой связи их экстраполированные значения при высокой энергии удивительно сближаются.}
\end{classroomqa}

\section{Бег к унификации}

Измерьте три калибровочные константы связи вблизи слабого масштаба и проследите их эволюцию, используя набор частиц некоторой предполагаемой теории. Поскольку каждая бета-функция логарифмична, огромное изменение энергии приводит лишь к умеренному изменению константы связи. Траектории Стандартной модели сближаются, но точно не встречаются. При низкоэнергетическом наборе частиц минимальной суперсимметрической Стандартной модели однопетлевые траектории сходятся гораздо ближе вблизи
\begin{equation}
  M_{\mathrm{GUT}}\sim10^{16}\,\mathrm{GeV}.
  \label{eq:l10-gut-scale}
\end{equation}

Этот график — свидетельство, а не доказательство. Пороговые эффекты от ещё не открытых частиц, поправки высших петель и детали объединённой теории смещают кривые. В лекции 2010 года подчёркивается сходимость примерно на процентном уровне для суперсимметрического спектра, потому что это делало суперсимметрию на слабом масштабе особенно привлекательной. Отсутствие суперпартнёров на коллайдерных энергиях теперь означает, что самая простая версия этого спектра не реализуется так, как первоначально ожидалось; точная унификация остаётся модельно-зависимой.

\begin{figure}[htbp]
\centering
\includegraphics[width=0.84\textwidth]{lecture_10_figure_10.png}
\caption{Качественный график констант связи на доске: сильная кривая падает, электрослабые кривые медленно эволюционируют к ней, а гравитация резко возрастает в далёком ультрафиолете.}
\lectureframe{01:16:29}
\label{fig:l10-unification-board}
\end{figure}

Почему считать окрестность точки встречи фундаментальной? Именно там исходные параметры могут стать наиболее простыми. Три не связанные между собой низкоэнергетические константы связи могли бы быть разными одетыми потомками одной высокоэнергетической константы связи, связанной с единственной калибровочной группой. Это численная подсказка, стоящая за более ранним вопросом об \(SU(5)\).

Гравитация появляется на той же логарифмической схеме как кривая, растущая гораздо быстрее. Она не пересекает калибровочные кривые точно на GUT-масштабе, но становится величиной порядка единицы в пределах нескольких декад по энергии. На логарифмической оси такая близость наводит на мысль: \(10^{16}\), \(10^{17}\), \(10^{18}\) и \(10^{19}\,\mathrm{GeV}\) несравненно ближе друг к другу, чем любое из них к электрослабому масштабу.

\section{Почему гравитация усиливается на малых расстояниях}

\subsection{Гравитирует энергия, а не только масса покоя}

Ньютоновская сила между двумя медленно движущимися массами равна
\begin{equation}
  F_G=\frac{Gm_1m_2}{r^2}.
  \label{eq:l10-newton-force}
\end{equation}
Для релятивистской размерностной оценки замените каждую массу покоя соответствующей энергией, делённой на \(c^2\):
\begin{equation}
  F_G(r)\sim\frac{GE_1E_2}{c^4r^2}.
  \label{eq:l10-relativistic-gravity-force}
\end{equation}
Это не заменяет общую теорию относительности; напряжения, поток импульса и давление тоже являются источниками гравитации. Для рассматриваемого вопроса это правильная оценка по порядку величины.

Локализуйте каждую частицу в области размера \(r\). Принцип неопределённости требует
\begin{equation}
  p\sim\frac{\hbar}{r}.
  \label{eq:l10-localization-momentum}
\end{equation}
При достаточно малых \(r\) режим становится ультрарелятивистским, так что
\begin{equation}
  E\simeq pc\sim\frac{\hbar c}{r}.
  \label{eq:l10-localization-energy}
\end{equation}
Подстановка двух таких энергий в
\eqref{eq:l10-relativistic-gravity-force} даёт
\begin{equation}
  F_G(r)
  \sim
  \frac{G\hbar^2}{c^2r^4}.
  \label{eq:l10-short-distance-gravity-force}
\end{equation}
Соответствующий масштаб потенциальной энергии равен
\begin{equation}
  U_G(r)\sim-\frac{G\hbar^2}{c^2r^3}.
  \label{eq:l10-short-distance-gravity-potential}
\end{equation}
Именно это различие было исправлено в ходе вычисления у доски: потенциал масштабируется как \(1/r^3\), тогда как сила — как \(1/r^4\).

\begin{classroomqa}
\audiencequestion{Для ультрарелятивистской частицы энергия равна
\(p^2/(2m)\)?\lecturetimestamp{01:22:18}}
\lecturerresponse{Нет. Это нерелятивистская формула. Из
\(E^2=p^2c^2+m^2c^4\) в пределе \(pc\gg mc^2\) получаем
\(E\simeq pc\). В сочетании с \(p\sim\hbar/r\) энергия локализации равна
\(\hbar c/r\).}
\end{classroomqa}

\begin{figure}[htbp]
\centering
\includegraphics[width=0.84\textwidth]{lecture_10_figure_11.png}
\caption{На доске релятивистские энергии подставляются в закон Ньютона с проверкой степеней \(c\) и \(r\).}
\lectureframe{01:23:29}
\label{fig:l10-gravity-derivation-board}
\end{figure}

\begin{classroomqa}
\audiencequestion{Сократились ли множители \(c\), и относился ли заявленный закон \(1/r^3\) к силе?\lecturetimestamp{01:23:38}}
\lecturerresponse{Расчёт на доске исправляется прямо по ходу дела. Две эффективные массы дают вклад \(E_1E_2/c^4\); при
\(E_i\sim\hbar c/r\) в знаменателе остаются две степени \(c\). Получающаяся сила равна \(G\hbar^2/(c^2r^4)\). Выражение \(1/r^3\) относится к потенциальной энергии.}
\end{classroomqa}

\begin{classroomqa}
\audiencequestion{Обсуждаемая гравитационная сила притягательна, верно?\lecturetimestamp{01:24:42}}
\lecturerresponse{Да. В уравнениях выше приведены модули. Ньютоновский гравитационный потенциал для положительных энергий содержит знак минус, так что в этой оценке сила притягательна.}
\end{classroomqa}

\begin{classroomqa}
\audiencequestion{Относительная сила гравитации возрастает потому, что локализация повышает энергию, тогда как сам электрический заряд не получает того же кинематического множителя?\lecturetimestamp{01:25:02}}
\lecturerresponse{Именно. Энергия — это гравитационный заряд. Сжатие квантового состояния увеличивает его импульс и энергию, а значит, увеличивает и то, что служит источником гравитации. Электрический заряд безразмерен и не приобретает непосредственно такого множителя \(1/r\), не считая его куда более медленной логарифмической ренормализации.}
\end{classroomqa}

\subsection{Планковское пересечение}

Электромагнитная сила между зарядами порядка единицы в естественных единицах имеет величину
\begin{equation}
  F_{\mathrm{EM}}(r)
  \sim
  \alpha\,\frac{\hbar c}{r^2}.
  \label{eq:l10-em-force-alpha}
\end{equation}
Приравнивая её к \eqref{eq:l10-short-distance-gravity-force}, получаем
\begin{equation}
  r^2\sim\frac{G\hbar}{\alpha c^3}.
  \label{eq:l10-gravity-em-crossing}
\end{equation}
В лекции безразмерный калибровочный множитель принимают порядка единицы, чтобы выделить фундаментальную гравитационную длину:
\begin{equation}
  \boxed{\ell_P=\sqrt{\frac{\hbar G}{c^3}}}
  \simeq1.616\times10^{-35}\,\mathrm{m}.
  \label{eq:l10-planck-length}
\end{equation}
Если сохранять реальное \(\alpha\), то точка равенства с электромагнетизмом сдвигается на логарифмической шкале на множитель порядка единицы, \(r\sim\ell_P/\sqrt\alpha\). В естественных единицах фундаментальное утверждение записывается чище:
\begin{equation}
\begin{aligned}
  \alpha_G(E)
  &\equiv\frac{GE^2}{\hbar c^5}
  =\left(\frac{E}{E_P}\right)^2,
  \\
  E_P
  &=\sqrt{\frac{\hbar c^5}{G}}
  \simeq1.22\times10^{19}\,\mathrm{GeV}.
\end{aligned}
  \label{eq:l10-gravitational-coupling}
\end{equation}
Гравитация становится сильной, когда \(E\) приближается к \(E_P\).

\begin{classroomqa}
\audiencequestion{В сравнении к единице приравнивали \(G\) или электрическую константу связи?\lecturetimestamp{01:27:02}}
\lecturerresponse{Электрическую константу связи рассматривали как число порядка единицы для грубой оценки. Постоянную Ньютона необходимо сохранить: это размерная величина, определяющая планковский масштаб.}
\end{classroomqa}

\begin{classroomqa}
\audiencequestion{Промежуточные степени \(r\), \(\hbar\) и \(c\) не выглядят размерностно согласованными. Как выглядит исправленный расчёт?\lecturetimestamp{01:27:47}}
\lecturerresponse{С самого начала используйте энергии. Локализация даёт
\(E\sim\hbar c/r\), откуда
\(U_G\sim GE^2/(c^4r)=G\hbar^2/(c^2r^3)\). Сравните это с
\(U_{\mathrm{EM}}\sim\alpha\hbar c/r\). Если взять
\(\alpha\) порядка единицы, получается
\(r^2=G\hbar/c^3\). И в аудитории, и на доске в итоге получают
\(r=\sqrt{G\hbar/c^3}\), планковскую длину.}
\end{classroomqa}

\begin{figure}[htbp]
\centering
\includegraphics[width=0.84\textwidth]{lecture_10_figure_12.png}
\caption{Исправленный итог размерностной проверки по ходу лекции:
\(r=\sqrt{\hbar G/c^3}\).}
\lectureframe{01:32:16}
\label{fig:l10-planck-length-board}
\end{figure}

\section{Проблема иерархии возвращается}

Калибровочные константы связи становятся сравнимыми вблизи очень высокого масштаба, а гравитация становится сильной вблизи планковского масштаба. Теперь эти огромные масштабы выглядят уже не столько как произвольные большие числа, сколько как области, где микроскопическое описание может упроститься. На их фоне электрослабое вакуумное значение \(F\) аномально мало:
\begin{equation}
  \frac{v}{E_P}\sim2\times10^{-17}
  \label{eq:l10-electroweak-planck-ratio}
\end{equation}
в единицах с \(c=1\). Использование GUT-масштаба вместо планковского меняет точное отношение, но не саму загадку.

Как только \(F\) мало, массы \(W\), \(Z\), заряженных лептонов и кварков тем самым определяются, поскольку калибровочная симметрия делает их массы пропорциональными ему. Частицы без такой защиты могут оставаться тяжелыми. Поэтому главный вопрос состоит не в том, почему каждая известная масса по отдельности мала. Он состоит в том, почему сектор Хиггса задает этот единственный малый масштаб нарушения симметрии.

\begin{classroomqa}
\audiencequestion{Если бы удалось вычислить \(F\), можно ли было бы тогда вычислить и все остальные массы?\lecturetimestamp{01:34:39}}
\lecturerresponse{Тогда были бы определены общий электрослабый масштаб и, при заданных калибровочных и юкавских константах связи, соответствующие массы частиц. Остается, однако, проблема: обычный ультрафиолетово-чувствительный расчет имеет тенденцию сдвигать параметр массы Хиггса к высокому фундаментальному масштабу, если только его не защищает симметрия или динамический механизм.}
\end{classroomqa}

Выражение «тонко настроенный» предвосхищает эту чувствительность. В описании с обрезанием скалярный параметр массы получает поправки грубого вида
\begin{equation}
  \delta\mu^2
  \sim
  \frac{C}{16\pi^2}\Lambda^2,
  \label{eq:l10-higgs-quadratic-sensitivity}
\end{equation}
где \(C\) объединяет калибровочные, юкавские и скалярные константы связи, а \(\Lambda\) обозначает масштаб, на котором эффективная теория заменяется. Если \(\Lambda\) близок к масштабу великого объединения или к масштабу Планка, получение \(\mu^2\sim v^2\), по-видимому, требует очень точной компенсации между исходным и радиационными вкладами. Это форма проблемы калибровочной иерархии с точки зрения естественности. Это редакторское завершение вычисления, обещанного в лекции на следующую четверть, а не утверждение о том, что артефакт конкретного регуляризатора сам по себе наблюдаем.

\begin{classroomqa}
\audiencequestion{Следует ли представлять длину Планка как конечный радиус экранирующего облака?\lecturetimestamp{01:35:03}}
\lecturerresponse{Нет. Диэлектрическая картина объясняет перенормировку калибровочного заряда, но аргумент Планка опирался на энергию локализации и размерную гравитационную константу связи. В этом выводе нет никакого буквального гравитационного экранирующего облака.}
\end{classroomqa}

\subsection{Заключительные вопросы о взаимодействиях}

В последние минуты лекция возвращается к потенциалу Хиггса и к тому, какие члены может содержать лагранжиан.

\begin{classroomqa}
\audiencequestion{Может ли потенциал Хиггса содержать \(\phi^6\) и при этом оставаться симметричным?\lecturetimestamp{01:35:20}}
\lecturerresponse{Да. Четные степени сохраняют игрушечную симметрию \(\phi\mapsto-\phi\), тогда как \(\phi^5\) нарушало бы ее. В четырех измерениях коэффициент при \(\phi^6\) имеет отрицательную массовую размерность, так что этот член подавлен высоким масштабом, например \(c_6\phi^6/\Lambda^2\). Это допустимое эффективное взаимодействие, хотя оно и не является перенормируемым по степенному счету.}
\end{classroomqa}

\begin{classroomqa}
\audiencequestion{Означают ли степени поля Хиггса процессы рождения и уничтожения, так что \(\phi^4\) предсказывает рассеяние Хиггс–Хиггс?\lecturetimestamp{01:35:52}}
\lecturerresponse{Да. После разложения около вакуума степени квантованного поля порождают вершины. Четвертичный член содержит четыреххиггсовскую вершину, а сдвинутый потенциал также порождает кубические взаимодействия. Знание их коэффициентов позволяет предсказывать процессы, такие как \(hh\to hh\), с учетом обычного диаграммного расчета.}
\end{classroomqa}

\begin{classroomqa}
\audiencequestion{Должен ли фундаментальный лагранжиан быть полиномиальным по своим полям, без степеней выше четвертой?\lecturetimestamp{01:36:09}}
\lecturerresponse{Для пертурбативно перенормируемой четырехмерной скалярной теории традиционный конечный список образуют операторы размерности не выше четырех. Современная теория эффективного поля менее строга: операторы большей размерности допустимы и кодируют физику малых расстояний, но их эффекты подавлены степенями высокого масштаба. Таким образом, редкое исключение существенно для старого правила.}
\end{classroomqa}

\begin{classroomqa}
\audiencequestion{Были ли члены самодействия опущены при обсуждении бегущих констант связи?\lecturetimestamp{01:36:44}}
\lecturerresponse{Они были учтены концептуально. Самодействия глюонов — именно то, что меняет знак по сравнению с КЭД и приводит к антиэкранированию. Самодействия Хиггса аналогичным образом вносят вклад в бег скалярных и других параметров Стандартной модели.}
\end{classroomqa}

\begin{classroomqa}
\audiencequestion{Есть ли перекрестные члены, в которых разные взаимодействия влияют на бег друг друга?\lecturetimestamp{01:37:23}}
\lecturerresponse{Да. На более высоком петлевом порядке бета-функции содержат смешанные произведения калибровочных, юкавских и скалярных констант связи. Теория возмущений остается полезной, пока эти безразмерные константы связи достаточно малы, поскольку последовательные петлевые порядки соответственно подавлены.}
\end{classroomqa}

\begin{classroomqa}
\audiencequestion{Что именно имелось в виду, когда говорилось, что \(F\) тонко настроено?\lecturetimestamp{01:37:44}}
\lecturerresponse{Полный ответ в лекции отложен. Проблема состоит в неустойчивости малого скалярного масштаба относительно гораздо больших ультрафиолетовых вкладов, что суммировано в \eqref{eq:l10-higgs-quadratic-sensitivity}. Суперсимметрия была одним из предложенных способов обеспечить компенсации, но сама проблема иерархии не выбирает подтвержденного решения.}
\end{classroomqa}

\begin{classroomqa}
\audiencequestion{Можно ли сейчас объяснить саму группу перенормировки?
\lecturetimestamp{01:37:56}}
\lecturerresponse{Не в оставшиеся минуты. Ее рабочее содержание уже проявилось: менять разрешение, давать константам связи течь и сохранять фиксированными физические предсказания. Полное изложение относится к следующему этапу курса.}
\end{classroomqa}

На этом эта последовательность, собственно, нас и оставляет. Механизм Хиггса сделал Стандартную модель внутренне непротиворечивой: он сохраняет хиральную калибровочную симметрию, придает массы слабым бозонам и фермионам и предсказывает физические взаимодействия Хиггса. Перенормировка затем учит нас тому, как эти взаимодействия меняются с масштабом, и намекает, что калибровочные силы могут иметь единое происхождение. Гравитация дает второй высокий масштаб посредством независимого аргумента. Ни один из этих успехов не объясняет, почему вакуумное значение Хиггса находится настолько ниже обоих. Поэтому заключительная лекция не закрывает тему; она указывает на то малое число, на котором должна сосредоточиться следующая тема.